\documentclass{article}

\usepackage[utf8]{inputenc}
\usepackage[turkish, shorthands = :!]{babel}
\usepackage[T1]{fontenc}
\usepackage{lmodern}

\usepackage{fancyhdr, lipsum, mathtools}

\usepackage{arydshln} % Draw dash-lines in array/tabular (ilaveli matrisleri arada kesikli çizgiyle yazabilmek için)

\renewcommand{\headrulewidth}{0 pt}
\pagestyle{fancy}
\fancyhf{}
\fancyhead[C]{\thepage}

\begin{document}

%%%%%%% SAYFA 1 %%%%%%%

\lipsum[1][1-3]
% burada asla boşluk olmamalı
\begin{align*}
4x - y + 2z		& = -1 \\
-x + (\ln 2)y - 3z	& = 0 % buraya \\ işareti konmamalı
\end{align*}
% eğer paragraf devam ediyorsa burada da boşluk olmamalı
\lipsum[1][4-5]

\bigskip
\hrule
\bigskip

\lipsum[1][1-3]
\[
\begin{array}{*{7}{r}} % array özetle varsayılanın matematik mod olduğu tabular
2x & + & 2y & + & 2z & = & 0 \\
-2x & + & 5y & + & 2z & = & 1 \\
8x & + & y & + & 4z & = & -1 % buraya \\ işareti koymak bir şey değiştirmiyor
\end{array}
\]
\lipsum[1][4-5]

\bigskip
\hrule
\bigskip

\lipsum[1][1-3]
\[
\arraycolsep = 0.66 mm
\begin{array}{*{7}{r}}
2x & + & 2y & + & 2z & = & 0 \\
-2x & + & 5y & + & 2z & = & 1 \\
8x & + & y & + & 4z & = & -1 \\
\end{array}
\]
\lipsum[1][4-5]

\newpage

%%%%%%% SAYFA 2 %%%%%%%

\lipsum[1][1-2]
$\begin{pmatrix}
a & b \\
c & d % buraya \\ işareti koymak bir şey değiştirmiyor
\end{pmatrix}$
\lipsum[1][3-5]

\bigskip
\hrule
\bigskip

\lipsum[1][1-3]
\[
\begin{pmatrix}
a & b \\
c & d % buraya \\ işareti koymak bir şey değiştirmiyor
\end{pmatrix}
\]
\lipsum[1][4-5]

\bigskip
\hrule
\bigskip

\[
\begin{pmatrix}
1 & -3 & 0 \\
-2 & 4 & 1 \\
5 & -2 & 2
\end{pmatrix}
\]

% array vasıtası ile matris girdileri güzelce hizalansa da parantezlerle girdiler arasında fazlaca boşluk oluşuyor (özellikle sol parantez ile ilk sütun arasında):
\[
\begin{pmatrix}
\begin{array}{rrr}
1 & -3 & \phantom{-} 0 \\ % \phantom üçüncü sütun biraz uzaklaşsın diye
-2 & 4 & 1 \\
5 & -2 & 2
\end{array}
\end{pmatrix}
\]

% mathtools paketinin sağladığı yıldızlı matris ortamlarını kullanmak hem daha iyi sonuç veriyor hem daha özlü:
\[
\begin{pmatrix*}[r]
%\begin{array}{rrr}
1 & -3 & \phantom{-}0 \\
-2 & 4 & 1 \\
5 & -2 & 2
%\end{array}
\end{pmatrix*}
\]

\bigskip
\hrule
\bigskip

\noindent
$2A^\top + C$ \\
$2A^\top \! + C$ \\ % tercihe şayan
$\det A^\top = \det A$

\bigskip
\hrule
\bigskip

\[
\left(
\begin{array}{c c c : c}
1 & -3 & 4 & 7 \\
0 & 1 & 2 & 2 \\
0 & 0 & 1 & 5 \\
\end{array}
\right)
\]

\end{document}